\scanpage{144}
\[
I^u:{\unsure{\mathcal A|A}}^G\longrightarrow\mathcal A^{G'} \tag*{(4)}
\]
On a aussi \deleted{On a aussi \unsure{$(2')$}}
\[
I^{1_G}=\deleted{\unsure{I}}\,1_{\mathcal A^G},\quad I^{uv}=I^uI^v \tag*{$(4')$}
\]
Les relations entre ces deux sortes de \deleted{opérations} \inserted{foncteurs} $\mathcal A^u,I^u$
et les \deleted{foncteurs} opérations $\otimes$ sont
\[
\left\{\begin{aligned}
\mathcal A^u(I^u(M))&=M^{[G':G]}\\
\deleted{\unsure{I^u(\mathcal A^u(M'))=M'{}^{[G':G]}}}&
\end{aligned}\right.\quad(M\in\mathcal A^G,\ G\subset G') \tag*{(5)}
\]
\deleted{\unsure{(6)}} On a d'ailleurs un hom fonctoriel
\deleted{\unsure{(6)} $M\longrightarrow$}
\[
M\longrightarrow\mathcal A^u I^u(M) \tag*{(6)}
\]
\[
M'\longrightarrow I^u\mathcal A^u(M') \tag*{(7)}
\]
\inserted{\unsure{$G$|2} hom. naturels !}

\margin{(8) Si $u$ est un isom, alors $I^u\unsure{=|\simeq}\mathcal A^{u^{-1}}$, donc $I^u$ et $\mathcal A^u$ sont des foncteurs inverses l'un de l'autre.}

\deleted{\unsure{Définissons|Définition}} (9) Soit $s$ la symétrie $G_1\times G_2\longrightarrow G_2\times G_1$. Alors $I^s(M\boxtimes M')=M'\boxtimes M$
% margin note written sideways
\margin{(10) Si $g\in G$, et $\sigma_g$ est l'automorphisme \uline{intérieur} défini par $g$, on a pour $M\in\mathcal A^G$ $M\simeq I^{\sigma_g}(M)\simeq\mathcal A^{\sigma_g}(M)$ L'isomorphisme \unsure{it.|itou|étant} est canonique et défini par $g$ opérant sur le vectoriel sous-jacent à $M$.}
% margin note written sideways
\margin{\unsure{formuler la propriété de cet isomorphisme}.}

Soit, pour des entiers $\geq1$, $p_1\ldots p_n$ de somme $p$
\[
u_{p_1,\ldots,p_n}:\unsure{\mathfrak S|\mathcal S}_{p_1}\times\ldots\times\mathfrak S_{p_n}\longrightarrow\mathfrak S_p \tag*{(\unsure{11|10})}
\]
l'hom (injectif) bien connu. Posons
pour $M\in\mathcal A^{\mathfrak S_p}$, $M'\in\mathcal A^{\mathfrak S_q}$
\[
M*M'=I^{u_{p,q}}(M\boxtimes M')\deleted{\unsure{\in\mathcal A^{\mathfrak S_{p+q}}}} \tag*{(\unsure{12|11}).}
\]
On a un isomorphisme canonique
\[
M*M'\simeq M'*M \tag*{(\unsure{13|12})}
\]
\deleted{canonique}, \deleted{\unsure{mais} à vrai dire \unsure{n'étant un tel} isomorphisme \unsure{qui est}} \inserted{défini par}
$\alpha$ : il existe $\alpha$ \deleted{des} unique de $\mathfrak S_{p+q}$
tel que $\alpha u_{p,q}\alpha^{-1}=u_{q,p}s_{p,q}$
-- où $s_{p,q}$ est la symétrie $\mathfrak S_p\times\mathfrak S_q\longrightarrow\mathfrak S_q\times\mathfrak S_p$.
% slanted
\margin{$*$ Noter aussi $\mathcal A^{u_{p,q}}(M*M')=\deleted{\unsure{M\boxtimes M'}}(M\boxtimes M')^{\frac{(p+q)!}{p!\,q!}}$ ce qui relie le $\boxtimes$ et le $*$, on peut \unsure{écrire} \deleted{$(M\otimes M')^{\unsure{\frac{(2p)!}{(p!)^2}}}$} $(M\odot M')^{\unsure{(2p)!/(p!)^2}}$ est défini par}

De façon analogue, on a un isomorphisme
canonique
\[
(M*M')*M''\simeq M*(M'*M'')\simeq I^{u_{p,q,r}}(M\boxtimes M'\boxtimes M'') \tag*{(14)}
\]
% archive sheet number: 143
% printed form (paper watermark visible in reverse)
\tikz[baseline=-.4ex,scale=.15]{\foreach \a in {0,45,...,315}{\draw (\a:.52) circle (.22);} \draw (0,0) circle (.21);} ARGENTINE \tikz[baseline=-.2ex,scale=.17]{\draw (0,0)--(1.6,0)--(1.8,.33)--(.5,.33)--(.32,.13)--(.92,.13);} TYPEWRITER \tikz[baseline=-.4ex,scale=.15]{\foreach \a in {0,45,...,315}{\draw (\a:.52) circle (.22);} \draw (0,0) circle (.21);} \unsure{A|AR}
